% ---
% titre: Problème de proportionnalité - la vente de pâtisserie
% theme: FA
% chapitre: Proportionnalité
% sous_chapitre: Résoudre des problèmes
% niveau: [10e]
% type: EVAL
% tags: [proportionnalite, p-question-TS]
% difficulte: 1
% corrige: true
% ---

\question
Voici une carte de la Suisse est à l'échelle 1:2\,400\,000. 

\begin{center}
\includegraphics[width=\linewidth]{Questions/Figures/CarteSuisse}
\end{center}


\begin{parts}

	\vspace{20pt}\part[5]
	Calcule la distance réelle, à vol d'oiseau, entre les villes de Lausanne et Zürich.
	
\vspace{10pt}{\footnotesize\raggedleft Espace pour ta démarche et tes calculs\par}
\begin{solutionorgrid}[140pt]
	
	\vspace{10pt}
	\begin{tabularx}{\linewidth}{|c*{2}{|>{\centering\arraybackslash}X}|}
	\hline\xrowht{20pt}
	\textbf{sur la carte [cm]}& 1 & 7.1 \\
	\hline\xrowht{20pt}
	\textbf{dans la réalité [cm]} & 2\,400\,000 & 17\,040\,000 \\
	\hline
	\end{tabularx}
	
	\vspace{10pt}
	7.1 $\cdot$ 2\,400\,000 = 17\,040\,000
	
	\vspace{10pt}
	Zürich se trouve à 17\,040\,000\,cm (170.4\,km) de Lausanne.
	
	\vspace{10pt}
	\begin{itemize}
	\item mesure de la distance en cm ($\pm$ 2\,mm) \textbf{(1pt)}
	\item tableau (avec nom de lignes) \textbf{(1pt)}
	\item calcul \textbf{(1pt)}
	\item réponse \textbf{(1pt)}
	\item phrase réponse \textbf{(1pt)}
	\end{itemize}
	\end{solutionorgrid}
	
	\newpage
	\part[8]
	Le train intercity 1 relie les villes de Lausanne, Fribourg, Bern et Zürich. Il met 50\,min pour relier Lausanne et Fribourg. Combien de temps dure le voyage Lausanne-Zürich considérant que le train roule à vitesse moyenne constante pendant tout le trajet?
	
	\vspace{10pt}
	{\small Indice: commence par calculer la distance entre Lausanne et Fribourg avant de calculer le temps de trajet entre Lausanne et Zürich.}
	
\vspace{10pt}{\footnotesize\raggedleft Espace pour ta démarche et tes calculs\par}
\begin{solutionorgrid}[140pt]
	
	\vspace{10pt}
	Distance Lausanne Fribourg: 
	
	\begin{tabularx}{\linewidth}{|c*{2}{|>{\centering\arraybackslash}X}|}
	\hline\xrowht{20pt}
	\textbf{sur la carte [cm]}& 1 & 2.1 \\
	\hline\xrowht{20pt}
	\textbf{dans la réalité [cm]} & 2\,400\,000 & 5\,040\,000 \\
	\hline
	\end{tabularx}
	
	\vspace{10pt}
	2.1 $\cdot$ 2\,400\,000 = 5\,040\,000
	
	\vspace{10pt}
	\begin{itemize}
	\item mesure de la distance en cm ($\pm$ 2\,mm) \textbf{(1pt)}
	\item tableau (avec nom de lignes) \textbf{(1pt)}
	\item calcul \textbf{(1pt)}
	\item réponse \textbf{(1pt)}
	\end{itemize}
	
	\vspace{10pt}
	Temps du trajet:
	
	\begin{tabularx}{\linewidth}{|c*{2}{|>{\centering\arraybackslash}X}|}
	\hline\xrowht{20pt}
	\textbf{durée [min]}& 50 & 169  \\
	\hline\xrowht{20pt}
	\textbf{distance [cm]} & 5\,040\,000 & 17\,040\,000 \\
	\hline
	\end{tabularx}
	
	\vspace{10pt}
	50 $\cdot$ 17\,040\,000 $\div$ 5\,040\,000 = 169
	
	\vspace{10pt}
	Le trajet Lausanne-Zürich dure 169\,min.
	
	\vspace{10pt}
	\begin{itemize}
	\item tableau (avec nom de lignes) \textbf{(1pt)}
	\item calcul \textbf{(1pt)}
	\item réponse \textbf{(1pt)}
	\item phrase-réponse \textbf{(1pt)}
	\end{itemize}
	\end{solutionorgrid}
\end{parts}

